\documentclass[letterpaper,12pt]{article}
\oddsidemargin -1.0cm \textwidth 16.5cm
%\usepackage[latin1]{inputenc}
%\usepackage[spanish]{babel}
\usepackage{tikz}
\usepackage[T1]{fontenc}
\usepackage[utf8]{inputenc}
\usepackage[spanish]{babel}
\usepackage{amsfonts,setspace}
\usepackage{amsmath}
\usepackage{amssymb, amsmath, amsthm}
\usepackage{tcolorbox}
\usepackage{comment}
%\usepackage[utf8]{inputenc}
\usepackage{amssymb}
\usepackage{anyfontsize}
\usepackage{bbm}
\usepackage{dsfont}
\usepackage{anysize}
\usepackage{multicol}
\usepackage{enumerate}
\usepackage{enumitem}
\usepackage{xcolor}
\usepackage{graphicx}
\usepackage{subfigure}
\usepackage{float}
\usepackage{fancyhdr}
\usepackage{algpseudocode}
\usepackage{algorithm}
\usepackage{listings}

\usepackage[font=small,labelfont=bf]{caption}
\usepackage[left=1.5cm,top=2cm,right=1.5cm, bottom=2cm]{geometry}
\usepackage{fancyhdr}
\pagestyle{fancy}

\definecolor{blue(pigment)}{rgb}{0.2, 0.2, 0.6}
\definecolor{codegreen}{rgb}{0,0.6,0}
\definecolor{codegray}{rgb}{0.5,0.5,0.5}
\definecolor{codepurple}{rgb}{0.58,0,0.82}
\definecolor{backcolor}{rgb}{0.95,0.95,0.92}
\definecolor{amber}{rgb}{1.0, 0.75, 0.0}
	\definecolor{ballblue}{rgb}{0.13, 0.67, 0.8}

\lstdefinestyle{mystyle}{
    backgroundcolor=\color{backcolor},
    commentstyle=\color{codegreen},
    keywordstyle=\color{magenta},
    numberstyle=\tiny\color{codegray},
    stringstyle=\color{codepurple},
    basicstyle=\ttfamily\footnotesize,
    breakatwhitespace=false,
    breaklines=true,
    captionpos=b,
    keepspaces=true,
    numbers=left,
    showspaces=false,
    showstringspaces=false,
    showtabs=false,
    tabsize=2
}

%Teoremas, Lemas, etc.
\theoremstyle{plain}
\newtheorem{teo}{Teorema}
\newtheorem{lem}{Lemma}
\newtheorem{prop}{Proposici\'on}
\newtheorem{cor}{Corolario}
\theoremstyle{definition}
\newtheorem{defi}{Definici\'on}
\newtheorem{eje}{Ejemplo}
\newtheorem{obs}{Observaci\'on}
\newcommand{\norm}[1]{\lVert#1\rVert}
\newcommand{\R}{\mathbb{R}}
\newcommand{\N}{\mathbb{N}}
\newcommand{\tiende}{\rightarrow}
\newcommand{\borel}{\mathcal{B}}
\newcommand{\proba}{\mathbb{P}}
\newcommand{\ds}{\displaystyle}
\newcommand{\partes}{\mathcal{P}}
\newcommand{\E}{\mathbb{E}}
\newcommand{\Var}{\mathbb{V}ar}
\newcommand{\1}{\mathbbm{1}}
% fin macros


%\fancypagestyle{plain}{%
%\fancyhf{}
%\lhead{\footnotesize\itshape\bfseries\rightmark}
%\rhead{\footnotesize\itshape\bfseries\leftmark}}
%\pagestyle{plain}
\usepackage{versions}
\includeversion{solution}

\let\epsilon\varepsilon

\begin{document}
\def\Pauta{1}
%%%%%ESCRIBIR 1 si quieren ver la Ayudantía con Pauta%%%%
%%%%%ESCRIBIR 0 si quieren ver sólo la Ayudantía%%%%

%Encabezado
\fancyhead[L]{\itshape{Escuela de Ingeniería}}
\fancyhead[R]{\itshape{Universidad de O'Higgins}}
\[
\vspace{-1.7cm}
\]

\begin{minipage}{11.5 cm}
\begin{flushleft}
\hspace*{-0.55cm}\textbf{ING1002-2  Cálculo Diferencial e Integral}\\
\hspace*{-0.7cm} \textbf{Profesor:} Gonzalo Flores G.\\
\hspace*{-0.7cm} \textbf{Ayudantes:} Juan I. Riquelme \& Cristian Acevedo R.\\
\end{flushleft}\end{minipage}

\begin{picture}(2,3)
\put(425,-8){\includegraphics[scale=0.23]{Imagenes/logo_uoh_2.png}}
\end{picture}

\vspace{-0.2cm}
\begin{center}
\LARGE\textbf{Ayudantía 4: Derivadas}\\
\large {\fontsize{13}{14}\selectfont 25 de Septiembre de 2025}
\end{center}
\vspace{-0.2cm}
\begin{enumerate}[label=\textbf{P\arabic*.}, itemsep=0cm]
\item Límites:
\begin{enumerate}
 \item $\displaystyle \lim_{x \to +\infty} \frac{e^{3x} - e^{-3x}}{e^{3x} + e^{-3x}}$
 
 \vspace{0.5em}
 \color{blue}
    \noindent{\bf Sol:}  Vamos a multiplicar el númerador y el denominador por $e^{-3x}$
$$
\lim_{x \to +\infty} \frac{e^{3x} - e^{-3x}}{e^{3x} + e^{-3x}} \cdot \left(\dfrac{e^{-3x}}{e^{-3x}}\right) = \lim_{x \to +\infty} \frac{1 - e^{-6x}}{1 + e^{-6x}}
 = \frac{1 - 0}{1 + 0} = 1
$$
\color{black}
\item
$\displaystyle \lim_{x \to +\infty}  \sqrt{3x^2 + 8x + 6} - \sqrt{3x^2 + 3x + 6} $

 \vspace{0.5em}
 \color{blue}
    \noindent{\bf Sol:} 
\[
= \lim_{x \to +\infty} \left( \sqrt{3x^2 + 8x + 6} - \sqrt{3x^2 + 3x + 6} \right) \cdot \frac{\sqrt{3x^2 + 8x + 6} + \sqrt{3x^2 + 3x + 6}}{\sqrt{3x^2 + 8x + 6} + \sqrt{3x^2 + 3x + 6}}  
\]
\[
= \lim_{x \to +\infty} \frac{5x}{\sqrt{3x^2 + 8x + 6} + \sqrt{3x^2 + 3x + 6}} = \lim_{x \to +\infty} \frac{5x}{\sqrt{3x^2 + 8x + 6} + \sqrt{3x^2 + 3x + 6}} \cdot \frac{\frac{1}{x}}{\frac{1}{x}}
\]
\[
= \lim_{x \to +\infty} \frac{5}{\sqrt{3 + \frac{8}{x} + \frac{6}{x^2}} + \sqrt{3 + \frac{3}{x} + \frac{6}{x^2}}}
= \frac{5}{2\sqrt{3}}
\]
\color{black}
\item \textbf{[Ejercicio Extra]}$\displaystyle \lim_{x \to +\infty} \frac{\sen^2(x)}{x^2 + 1}$

 \vspace{0.5em}
 \color{blue}
    \noindent{\bf Sol:} 
    Sabemos que \( 0 \leq \sen^2(x) \leq 1 \), entonces multiplicando por $\dfrac{1}{x^2+1}$2:
\[
0 \leq \frac{\sen^2(x)}{x^2 + 1} \leq \frac{1}{x^2 + 1}
\]
y como:
\[
\lim_{x \to +\infty} \frac{1}{x^2 + 1} = 0
\]
Por el \textbf{Teorema del Sandwich}, se concluye que:
\[
\lim_{x \to +\infty} \frac{\sen^2(x)}{x^2 + 1} = 0
\]    
\color{black}

\item $\displaystyle \lim_{x \to \frac{\pi}{2}} (1 + \cos x)^{3 \sec x}$

 \vspace{0.5em}
 \color{blue}
    \noindent{\bf Sol:}  Partamos recordando el siguiente \textbf{Límite notable:}
\[
\lim_{x \to \infty} \left(1 + \frac{1}{x}\right)^x = e
\]
Arreglemoslo un poquito, sea \(\frac{1}{x} = u\), cuando \( x \to \infty \), entonces \(u \to 0\), reemplazando:
\[
\Rightarrow \lim_{x \to \infty} \left(1 + \frac{1}{x}\right)^x = \lim_{u \to 0} (1 + u)^{\frac{1}{u}} = e
\]
Ahora resolvemos el límite, tenemos que
\[
\lim_{x \to \frac{\pi}{2}} (1 + \cos(x))^{3 \sec(x)}
= \lim_{x \to \frac{\pi}{2}} (1 + \cos(x))^{\frac{3}{\cos(x)}}
\]
Cambio de variable! Sea \(\cos(x) = u\), entonces cuando \(x \to \frac{\pi}{2} \Rightarrow u \to 0\), reemplazando:
\[
\Rightarrow \lim_{x \to \frac{\pi}{2}} (1 + \cos(x))^{\frac{3}{\cos(x)}} =\lim_{u \to 0} (1 + u)^{\frac{3}{u}} = \left( \lim_{u \to 0} (1 + u)^{\frac{1}{u}} \right)^3 = e^3 \quad \textcolor{blue}{\text{(Por el límite notable del inicio)}}
\]
\color{black}
\end{enumerate}
\item \textbf{[Ejercicio Extra]} Para cada una de las siguientes funciones, calcule todas las as\'{\i}ntotas (horizontales, verticales y oblicuas, si existen).
\begin{enumerate}
    \item $\displaystyle f(x) = \frac{x}{\sqrt{x^2 - 1}}$.
    
     \color{blue}
    \noindent{\bf Sol:} 

\textbf{Horizontal:} 
\[
 \lim_{x \to \infty} \frac{x}{\sqrt{x^2 - 1}} = 
\lim_{x \to \infty} \frac{\dfrac{x}{x}}{\sqrt{\dfrac{x^2}{x^2} - \dfrac{1}{x^2}}} =
\lim_{x \to \infty} \frac{1}{\sqrt{1 - \frac{1}{x^2}}}
= \frac{1}{\sqrt{1 - 0}} = 1
\]
Luego, tenemos asíntota horizontal en \( y = 1 \).

\textbf{Vertical:}
\[
x^2 - 1 = 0 \Rightarrow x = \pm 1 
\]
Luego, tenemos asíntota vertical en \( x = 1 \) y $x=-1$.

\textbf{Oblicua:} No hay, ya que $\displaystyle \lim_{x\to \infty} f(x)=1$, lo cual ya es la asíntota horizontal.
\color{black}
    
    \item $\displaystyle h(x) = e^{2/x}$.

 \color{blue}
    \noindent{\bf Sol:} 
    
\textbf{Horizontal:}
\[
\lim_{x \to \infty} e^{\frac{2}{x}} = e^0 = 1 \Rightarrow y = 1 \quad \text{es asíntota horizontal.}
\]
\textbf{Vertical:} \( x = 0 \), ya que 
\[
\lim_{x \to 0} e^{\frac{2}{x}} = e^{\infty} = \infty,
\quad \text{cumple la definición.}
\]
\textbf{Oblicua:} No hay, ya que $\displaystyle \lim_{x\to \infty} h(x)=1$, lo cual ya es la asíntota horizontal.
\color{black}
\end{enumerate}

\item Calcule la derivada de las siguientes funciones:
\begin{enumerate}
    \item \( f(x) = x^3 + 2e^x \sqrt{x} \)

 \color{blue}
    \noindent{\bf Sol:} Partamos con el primer término
    \[
    (x^3)' = 3x^2
    \]
    ahora para \( 2e^x \sqrt{x} = 2e^x x^{1/2} \), usamos la regla del producto:
    \[
    (2e^x x^{1/2})' = 2\left( (e^x)' \cdot x^{1/2} + e^x \cdot (x^{1/2})' \right)
    \]
    \[
    = 2\left( e^x \cdot x^{1/2} + e^x \cdot \frac{1}{2}x^{-1/2} \right)
    = 2e^x\left( x^{1/2} + \frac{1}{2}x^{-1/2} \right)
  \]
  \[ f'(x) = 3x^2 + 2e^x\left( x^{1/2} + \frac{1}{2}x^{-1/2} \right)
    \]

    
\color{black}

    \item \( f(x) = x(\sen x + \cos x) + 3e^x \)

 \color{blue}
    \noindent{\bf Sol:} Aplicamos la regla del producto al primer término
\[
\left[ x(\sen x + \cos x) \right]' =x'(\sen x + \cos x) + x(\sen x + \cos x)' = (\sen x + \cos x) + x(\cos x - \sen x) 
\]
derivamos el segundo término
\[
(3e^x)' = 3e^x
\]
\[
f'(x) = (\sen x + \cos x) + x(\cos x - \sen x) + 3e^x
\]
\color{black}

    \item \( f(x) = \dfrac{2}{6 + 6x + 3x^2} \)

 \color{blue}
    \noindent{\bf Sol:} Usamos la regla de la división:
\[
f'(x)= \frac{0 - 2(6 + 6x)}{(6 + 6x + 3x^2)^2} = \frac{-2(6 + 6x)}{(6 + 6x + 3x^2)^2}
\]
\color{black}

    \item \( f(x) = \dfrac{x^2 \cos(x)}{\sqrt{x}} \)

 \color{blue}
    \noindent{\bf Sol:} Primero vemos que
    
\[
(x^2 \cos(x))' = 2x \cos(x) - x^2 \sen(x)
\]
Además la derivada de $\sqrt{x}$, es\( = x^{1/2}\), entonces, aplicamos la regla de la división:
\[
f'(x) = \frac{(2x \cos(x) - x^2 \sen(x)) \cdot \sqrt{x} - x^2 \cos(x) \cdot \frac{1}{2}x^{-1/2}}{x}
\]
\color{black}

\end{enumerate}
\item Calcule las derivadas de las siguientes funciones usando Regla De La Cadena:
\begin{enumerate}
    \item\( f(x) = \sen^3(x) \cos^8(x) \)

 \color{blue}
    \noindent{\bf Sol:} 
Sea \( f(x) = u(x) v(x) \) con \( u(x) = \sen^3(x) \) y \( v(x) = \cos^8(x) \). Derivamos usando la regla del producto y de la cadena:
\[
u'(x) = 3 \sen^2(x) \cos(x)
\]
\[
v'(x) = 8 \cos^7(x)(-\sen x) = -8 \cos^7(x) \sen(x)
\]
Por lo tanto,
\[
f'(x) = u'(x)v(x) + u(x)v'(x) = 3 \sen^2(x) \cos(x) \cdot \cos^8(x) - 8 \sen^3(x) \cos^7(x) \sen(x)
\]
\[
 = 3 \sen^2(x) \cos^9(x) - 8 \sen^4(x) \cos^7(x)
\]
\color{black}
    
    \item  \( g(x) = \dfrac{\ln\left( \sqrt[5]{x^2 + 1} \right)}{x^n} \)

 \color{blue}
    \noindent{\bf Sol:} 

Primero simplificamos el logaritmo:
\[
\ln\left( \sqrt[5]{x^2 + 1} \right) = \frac{1}{5} \ln(x^2 + 1)
\]
Entonces
\[
g(x) = \frac{1}{5} \cdot \frac{\ln(x^2 + 1)}{x^n}
\]
Usamos la regla de la división con
\[
u(x) = \ln(x^2 + 1), \quad v(x) = x^n
\]
Calculamos las derivadas:
\[
u'(x) = \frac{2x}{x^2 + 1}, \quad v'(x) = n x^{n-1}
\]
Derivamos usando la regla del cociente:
\[
g'(x) = \frac{1}{5} \cdot \frac{u'(x)v(x) - u(x)v'(x)}{[v(x)]^2} = \frac{1}{5} \cdot \frac{\frac{2x}{x^2 + 1} \cdot x^n - \ln(x^2 + 1) \cdot n x^{n-1}}{x^{2n}}
\]  
\color{black}

\item $u(x)=2^x,\ \text{Hint: Use que $a=e^{\ln(a)},\forall a>0$}.$

 \color{blue}
    \noindent{\bf Sol:} 
\begin{align*}
u(x) &= 2^x \\
     &= \left(e^{\ln(2)}\right)^x \quad \text{(ya que $2 = e^{\ln(2)}$)} \\
     &= e^{x \ln(2)}
\end{align*} 
Luego 
\begin{align*}
u(x) &= e^{x \ln(2)} \\
u'(x) &= \frac{d}{dx} \left( e^{x \ln(2)} \right) \\
      &= \ln(2) \cdot e^{x \ln(2)} \\
      &= \ln(2) \cdot 2^x
\end{align*}
\color{black}


    
    \item \( h(x) = \ln(1 + e^{-3x + 2}) \)

 \color{blue}
    \noindent{\bf Sol:} 
Sea \( u(x) = 1 + e^{-3x + 2} \), entonces
\[
h(x) = \ln(u(x)) \implies (h(x))'=\frac{1}{u'(x)}
\]
Calculamos \( u'(x) \):
\[
u'(x) = \left(1 + e^{-3x + 2}\right)' = -3 e^{-3x + 2}
\]
Por lo tanto,
\[
h'(x) = \frac{-3 e^{-3x + 2}}{1 + e^{-3x + 2}}
\]

\color{black}
    
\end{enumerate}
\item Encontrar la ecuaci\'on de la recta tangente a la curva $x^2 + y^2 - xy = 1$ en el punto $(1,1)$.

 \color{blue}
    \noindent{\bf Sol:} 
Diferenciamos implícitamente:
\[
2x + 2y \frac{dy}{dx} - \left( y + x \frac{dy}{dx} \right) = 0
\]
\[
2y \frac{dy}{dx} - x \frac{dy}{dx} = y - 2x
\]
Sacamos factor común \(\dfrac{dy}{dx}\):
\[
\frac{dy}{dx}(2y - x) = y - 2x
\]
\[
\frac{dy}{dx} = \frac{y - 2x}{2y - x}
\]
Evaluamos en \((1,1)\):
\[
\left.\frac{dy}{dx}\right|_{(1,1)} = \frac{1 - 2(1)}{2(1) - 1} = \frac{1 - 2}{2 - 1} = \frac{-1}{1} = -1
\]

Ecuación de la recta tangente:
\[
y - 1 = -1 (x - 1) \implies y = -x + 2
\]
\color{black}
\item Encuentre la ecuaci\'on de la recta tangente a la gr\'afica de la funci\'on $f(x) = \arctan x$ en $x = 1$.

 \color{blue}
    \noindent{\bf Sol:} 
Partamos con la derivada de $\arctan x$
\[
(\arctan x)' = \frac{1}{1 + x^2}
\]

Evaluamos en \( x = 1 \):
\[
f'(1) = \frac{1}{1 + 1^2} = \frac{1}{2}
\]

Calculamos el punto para $y$:
\[
f(1) = \arctan(1) = \frac{\pi}{4}
\]

Ecuación de la recta tangente:
\[
y - \frac{\pi}{4} = \frac{1}{2}(x - 1) \implies y = \frac{1}{2}x + \frac{\pi}{4} - \frac{1}{2}
\]

\end{enumerate}

\end{document}
